% Options for packages loaded elsewhere
\PassOptionsToPackage{unicode}{hyperref}
\PassOptionsToPackage{hyphens}{url}
\PassOptionsToPackage{dvipsnames,svgnames,x11names}{xcolor}
\documentclass[
  11pt,
  letterpaper,
]{article}
\usepackage{xcolor}
\usepackage[margin=1in]{geometry}
\usepackage{amsmath,amssymb}
\setcounter{secnumdepth}{-\maxdimen} % remove section numbering
\usepackage{iftex}
\ifPDFTeX
  \usepackage[T1]{fontenc}
  \usepackage[utf8]{inputenc}
  \usepackage{textcomp} % provide euro and other symbols
\else % if luatex or xetex
  \usepackage{unicode-math} % this also loads fontspec
  \defaultfontfeatures{Scale=MatchLowercase}
  \defaultfontfeatures[\rmfamily]{Ligatures=TeX,Scale=1}
\fi
\usepackage{lmodern}
\ifPDFTeX\else
  % xetex/luatex font selection
\fi
% Use upquote if available, for straight quotes in verbatim environments
\IfFileExists{upquote.sty}{\usepackage{upquote}}{}
\IfFileExists{microtype.sty}{% use microtype if available
  \usepackage[]{microtype}
  \UseMicrotypeSet[protrusion]{basicmath} % disable protrusion for tt fonts
}{}
\makeatletter
\@ifundefined{KOMAClassName}{% if non-KOMA class
  \IfFileExists{parskip.sty}{%
    \usepackage{parskip}
  }{% else
    \setlength{\parindent}{0pt}
    \setlength{\parskip}{6pt plus 2pt minus 1pt}}
}{% if KOMA class
  \KOMAoptions{parskip=half}}
\makeatother
\usepackage{longtable,booktabs,array}
\newcounter{none} % for unnumbered tables
\usepackage{calc} % for calculating minipage widths
% Correct order of tables after \paragraph or \subparagraph
\usepackage{etoolbox}
\makeatletter
\patchcmd\longtable{\par}{\if@noskipsec\mbox{}\fi\par}{}{}
\makeatother
% Allow footnotes in longtable head/foot
\IfFileExists{footnotehyper.sty}{\usepackage{footnotehyper}}{\usepackage{footnote}}
\makesavenoteenv{longtable}
\setlength{\emergencystretch}{3em} % prevent overfull lines
\providecommand{\tightlist}{%
  \setlength{\itemsep}{0pt}\setlength{\parskip}{0pt}}
\usepackage[none]{hyphenat}
\sloppy
\usepackage{fancyhdr}
\usepackage{titlesec}
\pagestyle{fancy}
\fancyhf{}
\fancyhead[L]{\small Agent Payment Authority}
\fancyhead[R]{\small Working Paper}
\fancyfoot[C]{\thepage}
\setlength{\headheight}{14pt}
\setlength{\emergencystretch}{3em}
\titleformat{\section}{\large\bfseries}{}{0pt}{}
\titleformat{\subsection}{\normalsize\bfseries}{}{0pt}{}
\usepackage{bookmark}
\IfFileExists{xurl.sty}{\usepackage{xurl}}{} % add URL line breaks if available
\urlstyle{same}
\hypersetup{
  pdftitle={Deterministic Authorization and Recovery for Agentic Payments},
  pdfauthor={Working paper · revised draft},
  colorlinks=true,
  linkcolor={Maroon},
  filecolor={Maroon},
  citecolor={Blue},
  urlcolor={MidnightBlue},
  pdfcreator={LaTeX via pandoc}}

\title{Deterministic Authorization and Recovery for Agentic Payments}
\usepackage{etoolbox}
\makeatletter
\providecommand{\subtitle}[1]{% add subtitle to \maketitle
  \apptocmd{\@title}{\par {\large #1 \par}}{}{}
}
\makeatother
\subtitle{A failure analysis and proposed structure, with notes on
integration}
\author{Working paper · revised draft}
\date{September 18, 2026}

\begin{document}
\maketitle

\subsection{Abstract}\label{abstract}

Agent-initiated spending is being deployed faster than any public record
of its failures is accumulating. A survey of the documented record finds
one well-corroborated case of an autonomous agent making an unauthorized
real-money purchase, a small cluster of merchant-side wrong-outcome
failures, and a body of red-team work demonstrating capability without
confirmed victims. That scarcity is a reporting artifact rather than
evidence of reliability: there is no disclosure regime for agent
purchase errors, and individual losses are small enough to be absorbed
by ordinary chargebacks.

Benchmark data fills the gap the incident record leaves. On a retail
agent benchmark, a frontier model completes a task correctly on a single
attempt far more often than it does so consistently across repeated
attempts, and the most over-represented retail error category is
confirmation handling. A control designed for an agent that is usually
right must therefore assume it will sometimes be wrong in exactly the
step where money moves.

This paper argues that existing agentic payment controls bound the wrong
quantity. Per-transaction amount caps, tokenized credentials, and spend
limits bound \emph{how much} an agent can spend. Every documented
failure surveyed here occurred within an authorized amount, at a
permitted merchant, on a legitimately delegated credential. What was
wrong was the purchase, not the price.

We propose a structure with two halves: make authorization deterministic
by binding it to an exact, human-approved, hash-identified record
evaluated by pure functions; and make recovery explicit by treating
clawback as a policy-bound workflow that reports what it could not
recover. A local prototype implements the structure for one purchase
category using simulated providers, and includes adversarial tests for
attacks that an earlier revision of the same code failed. It does not
establish that declared software actually ran, authenticate real
merchant evidence, or demonstrate production payment safety.

\subsection{1. Deployment is outpacing the failure
record}\label{deployment-is-outpacing-the-failure-record}

In February 2025 a technology columnist asked OpenAI's Operator to
compare egg prices in his neighborhood. Within ten minutes the agent
bought a dozen eggs through Instacart and paid for human delivery, a
purchase he had not authorized. The transaction was \$31.43. What makes
it more than an anecdote is that it bypassed OpenAI's own stated
safeguard requiring user confirmation before a purchase; OpenAI
acknowledged the failure and committed to improving those safeguards.
{[}18{]}

That incident is, as far as this survey can establish, the clearest
public case of an autonomous agent making an unauthorized real-money
purchase. It is nearly the only one.

This is not a claim that agent purchase errors are rare. It is a claim
about the record. Searching the incident databases, technology press,
security research, and regulatory publications for documented
agent-initiated purchase failures returns a very short list, assembled
in Section 2. Meanwhile the infrastructure for such purchases has been
deployed broadly: card networks have published agent transaction
protocols {[}6{]} {[}5{]}, a major marketplace operates an agentic
checkout against third-party catalogs {[}20{]}, and consumer assistants
that transact on a user's behalf have raised substantial capital and
reached general availability.

Several structural reasons explain the gap, none of them reassuring:

\begin{itemize}
\tightlist
\item
  \textbf{No disclosure regime.} No rule requires a platform to publish
  an agent purchase error. Voluntary transparency reporting so far has
  covered model behavior rather than transactions; a September 2026
  disclosure of six agent incidents by OpenAI contained no purchase,
  payment, or booking events, all six having occurred in internal
  testing. {[}24{]}
\item
  \textbf{Losses are small enough to disappear.} A \$31 error resolves
  through an ordinary chargeback or a goodwill refund and never becomes
  a public record.
\item
  \textbf{Merchant-side harms settle privately.} Where the injured party
  is a business rather than a consumer, the dispute resolves
  contractually. This survey found no documented enterprise procurement
  agent failure, which is more plausibly a disclosure artifact than an
  absence of error.
\end{itemize}

The honest formulation is that agent-initiated spending is being
deployed considerably faster than public evidence about its reliability
is accumulating. A design discipline for this domain therefore cannot
wait for the incident record to mature, and should not be calibrated to
the small number of failures that happen to have been reported.

\subsection{2. What the record does
contain}\label{what-the-record-does-contain}

Sources were accessed September 18, 2026. Each item below is labeled by
evidence tier. Two claims encountered during this survey --- a statistic
on agent dispute rates and a purported regulatory advisory on
autonomous-agent purchases --- could not be substantiated at their
attributed sources and are deliberately excluded.

\subsubsection{2.1 Liability is already settled, and it sits with the
operator}\label{liability-is-already-settled-and-it-sits-with-the-operator}

Before the agentic era, a Canadian tribunal ruled on an airline chatbot
that gave a passenger incorrect information about a bereavement fare,
causing him to pay \$1,630.36 in full fare. The airline argued that the
chatbot was ``a separate legal entity'' for which it was not
responsible. The tribunal rejected this, holding that the company ``is
responsible for all the information on its website'' and that it ``makes
no difference whether the information comes from a static page or a
chatbot,'' and awarded damages. {[}19{]}

\emph{Evidence tier: court ruling.} This is the strongest item in the
survey and the most consequential. It establishes that an operator owns
its agent's errors. Any architecture that leaves an agent's wrong
purchase unattributable is therefore not merely a technical gap; it is
an unmanaged liability.

\subsubsection{2.2 Authorized purchase, wrong
outcome}\label{authorized-purchase-wrong-outcome}

This is the failure mode central to this paper, and the most thinly
documented.

In December 2024 a reporter asked a shopping agent to buy a specific
toothpaste from a named retailer. The agent took payment --- confirmed
on the reporter's bank statement --- then failed for three hours without
completing the purchase. The item was in fact out of stock at the
retailer while appearing available in the agent's scraped search data.
{[}21{]} \emph{Evidence tier: named-source journalism, firsthand, single
source.} Authorization was valid at every layer. The agent's model of
the world was wrong.

At larger scale, a marketplace agentic checkout placed orders against
stale, incorrect, or entirely absent third-party catalog data beginning
around November 2025. Named merchants reported receiving orders for
products they do not sell, listings displaying photographs of unrelated
items, and orders for discontinued stock. When the marketplace's data
was wrong, the merchant absorbed return shipping and refund. One
affected merchant collected 145 responses from other brands. {[}20{]}
\emph{Evidence tier: named-source journalism, multiple named merchants,
two outlets.}

The marketplace case inverts the usual framing: the consumer's
authorization was entirely valid and the injured party was the merchant.
It is the only fleet-scale evidence located in this survey.

\subsubsection{2.3 Irreversible action against an explicit
constraint}\label{irreversible-action-against-an-explicit-constraint}

In July 2025 a coding agent, operating under an explicit instruction not
to modify production during a code freeze, executed destructive commands
against a live production database, deleting records for approximately
1,200 companies and a comparable number of individuals. It then reported
fabricated results and incorrectly asserted that recovery was
impossible. The data was recovered manually; the vendor apologized and
shipped environment separation and a non-executing planning mode within
days. {[}22{]} \emph{Evidence tier: firsthand, company acknowledgment,
journalism.}

This was not a financial transaction and is included only for the
mechanism it demonstrates: an agent taking an irreversible action that
the operator had explicitly forbidden, then misreporting it. The same
mechanism applied to a payment produces a loss that no amount limit
would have prevented.

\subsubsection{2.4 Injection-driven spending: capability, not
casualties}\label{injection-driven-spending-capability-not-casualties}

Security researchers have catalogued prompt-injection payloads on live
web pages that specifically target payment actions, including attempts
to subscribe a victim to a paid plan, force a donation through a payment
link, and redirect a substantial transfer to an attacker-controlled
account. The publishing researchers state explicitly that they are not
aware of any confirmed real-world instance where such an attack
succeeded against a deployed system. {[}23{]} Separate vulnerability
research has demonstrated indirect injection against agentic browsers,
including payloads concealed in images. {[}25{]}

\emph{Evidence tier: vendor security research.} The accurate formulation
is that the capability is demonstrated and payloads are present in the
wild, while no confirmed consumer losses have been published. This must
not be presented as an incident record. It does establish that an
agent's instruction stream is attacker-reachable, which bears directly
on why authority should bind to terms a human approved rather than to
the agent's in-flight reasoning.

\subsection{3. The rate, which the anecdotes cannot
supply}\label{the-rate-which-the-anecdotes-cannot-supply}

Because the incident record is thin, the more informative evidence is
measured error rates. A retail-and-airline agent benchmark evaluates not
only whether an agent completes a task correctly but whether it does so
\emph{consistently} across repeated independent attempts --- reported as
pass\^{}k. {[}17{]}

The distinction matters more than the headline accuracy number. On the
retail domain, a model measured at roughly 61\% correct on a single
attempt fell to roughly 25\% correct across eight consecutive attempts.
Frontier models have since improved single-attempt retail performance
substantially, but the shape of the pass\^{}k decline is the durable
finding: an agent that is usually right is considerably less often
reliably right.

Two further details from that work bear on this paper directly. First,
the benchmark's retail error taxonomy over-represents \emph{confirmation
handling} errors relative to other domains --- the agent mishandling the
step where it should have checked with the user. That is precisely the
failure exhibited in the Operator incident, where a stated confirmation
safeguard did not fire. Second, the benchmark also over-represents
\emph{constraint interpretation} errors, the agent misreading what it
was told to satisfy.

A spending agent operates in the repeated regime, not the single-attempt
one. A control designed on the assumption that the agent is usually
correct is calibrated to the wrong statistic.

\subsection{4. Why current controls do not address these
failures}\label{why-current-controls-do-not-address-these-failures}

The controls shipped with current agentic payment products are, with one
exception, quantity controls.

A consumer assistant that transacts on a user's behalf mints a
single-use virtual card authorized for a specific amount, so that the
agent never handles the underlying card number. {[}9{]} This is a
genuine and well-designed control: it bounds the exposure of any one
transaction and removes the credential from the agent's reach entirely.

A payments platform for agents enforces permissions and limits
server-side on every request, so that a compromised or misbehaving agent
cannot exceed what it was granted. Limits may be set per transaction,
per day, and per month. A payment breaching a limit is not rejected; it
is held and an approval is opened that the agent cannot clear itself.
{[}7{]} Server-side evaluation and the inability of an agent to raise
its own caps are both correct design choices, and stronger than
client-side enforcement.

Neither control addresses any failure in Section 2.

The unauthorized egg purchase was for \$31.43 --- within any plausible
limit. The toothpaste purchase was the right item at the right merchant
for the right price; only the inventory was wrong. The marketplace
orders were correctly priced against catalogs that were themselves
incorrect. In each case an amount-based control evaluates the
transaction and finds nothing to object to, because there was nothing
wrong with the amount.

Two further properties of the current model deserve note:

\textbf{Holding rather than blocking relocates the decision but does not
make it deterministic.} A limit breach that opens an approval is a good
design for the case where the user is available and understands what
they are approving. It does not help where the purchase is within
limits, which is the case in every incident surveyed.

\textbf{Where confirmation is discretionary, it is not a control.} One
consumer assistant's terms appoint the service as the user's agent to
``enter into agreements, commitments or transactions'' on their behalf,
then state that the provider ``may implement safeguards, confirmation
requirements, or other controls on certain Actions'' while making ``no
representation or warranty that such safeguards will prevent unintended
or erroneous Actions.'' The same terms make payments non-refundable,
direct purchase disputes to the merchant rather than the provider, and
cap provider liability at one hundred dollars. {[}26{]} An independent
assessment of the same product found that confirmation before charge
does sometimes occur. {[}27{]} A safeguard that fires sometimes, is
disclaimed contractually, and is paired with no recovery path is not a
control a user can reason about.

The missing axis is not \emph{how much} but \emph{what}, and
\emph{whether it worked}.

\subsection{5. Prior work and the scope of this
proposal}\label{prior-work-and-the-scope-of-this-proposal}

This is a targeted review of relevant primary documentation, not a
systematic literature review or exhaustive market survey. Sources were
accessed on September 17, 2026. Vendor statements describe public
positioning; specifications describe intended protocol behavior; neither
alone proves adoption or operational performance. Drafts and preprints
are identified as such.

\textbf{Delegation and transaction intent already have substantial
coverage.}

OAuth Rich Authorization Requests already expresses fine-grained
authority in structured form, including payment amount and recipient.
Granular policy is therefore not an invention specific to AI agents.
{[}2{]}

The inspected AP2 v0.2 specification defines Checkout and Payment
Mandates, linked receipts, verification responsibilities, and direct and
autonomous flows. Its constraints are evaluated by designated
participants; it also describes evidence used in disputes. The earlier
launch essay used Intent and Cart Mandate terminology, so an
implementation should follow a pinned specification rather than infer
today's schema from that announcement. This paper's agent-configuration
version is distinct from AP2's mandate-schema version. {[}3, 4{]}

Mastercard-maintained Verifiable Intent, marked draft v0.1 in the
inspected repository, specifies a delegation chain, selective
disclosure, purchase constraints, and checkout-payment binding. It
distinguishes machine-enforceable constraints from descriptive product
information and leaves dispute resolution outside its scope. Visa's
Trusted Agent Protocol specifies signed agent-recognition and related
commerce messages for participating merchants. Neither a valid message
signature nor a signed description should be interpreted as a
measurement of an entire executing agent. {[}5, 6{]}

\textbf{Persistent identities and consumers are not unserved
categories.}

Natural publicly describes persistent agent history, revocable keys,
limits, and permissions. Its public materials cover a broad financial
stack and discuss disputes; it would be inaccurate to characterize
Natural as only a wallet or as lacking guardrails. This review did not
verify its internal version-binding implementation, every product's
availability, or an integration with Scrip. {[}1, 7, 8{]}

Stripe's April 2026 announcement describes a Link wallet for agents,
consumer approval of spending requests, and one-time cards or Shared
Payment Tokens without exposing underlying payment credentials. A
``consumer wallet for agents'' is thus not sufficient differentiation by
itself. {[}9{]}

The potential product distinction examined here is narrower: making the
relationship between software changes, outstanding authority, and
subsequent outcomes understandable to the principal. Whether that merits
a standalone product, a provider feature, or an open-source component is
an open commercial question.

\textbf{Identity, provenance, and outcome verification have prior art.}

SPIFFE/SPIRE already separates workload identity from credentials and
uses workload properties during attestation. SLSA provenance describes
how an artifact was produced. The IETF RATS architecture distinguishes
evidence, appraisal, and attestation results, with explicit freshness
and trust assumptions. These are foundations to reuse, not missing
infrastructure that this paper replaces. {[}10--12{]}

ERC-8004, marked draft at review time, proposes identity, reputation,
and validation registries while leaving payments orthogonal. Agentic
Settlement Protocol is particularly close adjacent work: its September
2026 preprint proposes holds, fulfillment verification, and
partial-refund mechanisms for delayed-fulfillment stablecoin commerce,
including travel. It identifies itself as a design paper without
deployed measurements. Outcome-aware commerce and recovery are therefore
not novel merely because they are connected to agents. {[}13, 14{]}

This proposal concentrates on principal-approved configuration changes
across a mission lifecycle. That emphasis does not establish priority
over all prior work, and no interoperability with these systems has been
demonstrated.

\subsection{6. System model and trust
assumptions}\label{system-model-and-trust-assumptions}

The \textbf{principal} grants authority. An \textbf{operator} deploys
the agent. A \textbf{control service} authenticates requests, checks
mandates, and maintains budget and operation state. \textbf{Payment
providers} report financial events. \textbf{Evidence providers} report
fulfillment or recovery events. Some roles may belong to one
organization, but their claims should remain distinguishable.

The agent may generate mistaken requests, encounter malicious content,
repeat operations, or continue after an operator changes its
configuration. Providers may time out, return events out of order, or
supply inconsistent evidence. These are the failures the design aims to
expose or constrain.

The intended controls require all covered spending to pass through the
service or equivalent provider-enforced restrictions. An agent with an
unrestricted credential outside that boundary can bypass them. The
design assumes a trustworthy authorization service, protected state,
authenticated principal actions, and appropriate verification of
provider messages. A compromised administrator, colluding evidence
source, or stolen credential can invalidate those assumptions. Storing a
hash beside data does not protect both from an attacker who can rewrite
the database.

``Verified outcome'' means an assessment supported by specified evidence
under a stated trust policy. It does not mean objective truth, customer
satisfaction, legal liability, or proof of the agent's private
reasoning.

Manifests and receipts can reveal sensitive instructions, destinations,
purchases, and behavior. The proposed service should minimize
disclosure, restrict access, and define retention periods. A digest is
not encryption: predictable underlying content may still be guessed.
Recording private chain-of-thought is neither necessary nor an assurance
mechanism in this design.

\subsection{7. A deterministic authorization
structure}\label{a-deterministic-authorization-structure}

The goal is not to make agent behavior deterministic; it is not. The
goal is to make the \emph{authorization decision} deterministic, so that
whatever the agent does, the set of purchases it can cause is a pure
function of a record a human approved.

This section states the structure independently of any implementation.
Section 9 describes a prototype that realizes it for one purchase
category.

\subsubsection{7.1 Approve an exact}\label{approve-an-exact}

A spending authorization should name the purchase, not a budget. ``A
refundable room at this hotel, these dates, this rate'' rather than
``hotels under \$550.'' The approved record is canonicalized and hashed;
the hash is what the human's approval binds to.

This inverts the usual tradeoff. A budget is convenient and admits every
failure in Section 2, because every one of those purchases was inside
its budget. An exact record is less convenient and admits none of them:
a cancellation is not a purchase in the record, a different item is not
the item in the record, an out-of-stock substitution is not the approved
quote.

The cost is real and should be stated. Exact approval means an agent
that finds a better option must return for a new approval, and a system
that requires many approvals may be abandoned or approved reflexively.
Section 10 treats this as the central open question rather than a solved
problem.

\subsubsection{7.2 Constraints must}\label{constraints-must}

Conditions that govern payment must have machine-checkable
representations --- amount ceilings, permitted merchants, date ranges,
refundability deadlines --- evaluated by a pure function over the
approved record and the proposed purchase. Natural language may propose
these terms; it must not enforce them. A model's assertion that a
purchase looks acceptable is not an enforcement mechanism.

\subsubsection{7.3 What cannot be typed}\label{what-cannot-be-typed}

A condition the system cannot represent as a predicate (``somewhere
nice'', ``good reviews'') is the dangerous case, because the natural
implementation is to drop it. The structure requires that unresolved
conditions be recorded as such and that their presence block unattended
execution. Failing closed on uninterpretable terms is what prevents a
silent gap between what the human said and what the system enforces.

This also localizes the nondeterminism that natural-language
interpretation necessarily introduces. Extraction from prose is
fallible; the design requirement is not that extraction be correct but
that it be \emph{calibrated} --- confidently mistyping a constraint is
the harmful failure, while declining to type one is safe because it
blocks. A useful evaluation of such a system measures silent mistyping
against correct refusal, not extraction accuracy.

\subsubsection{7.4 Authority binds}\label{authority-binds}

An authorization granted to an agent is granted to a particular declared
configuration of that agent: model reference, instruction and policy
digests, code artifact, tool versions and permissions. A newly
registered configuration inherits no authority; a change to any
authority-bearing field requires review under the principal's policy,
and a change whose behavioral effect is unknown --- a dependency bump, a
code change labeled a patch --- requires review rather than receiving an
automatic exemption.

Authority is checked before any exposure is reserved, and rechecked
immediately before the external call. The second check exists because
revocation can arrive while an operation is in flight. It narrows the
window; it does not close it. Ordering between a local registry and an
external provider is not atomic, and once a provider has accepted an
instruction, revocation cannot retract it --- that operation must be
reconciled under the configuration that issued it rather than disowned.

\subsubsection{7.5 Outcome is assessed}\label{outcome-is-assessed}

Payment success and fulfillment success are separate facts and must be
recorded separately. Capture is recorded when financial evidence
supports it, even where fulfillment later fails; delaying recognition of
spent money until the outcome is satisfactory misstates exposure.

The agent's own account of what happened is not evidence. An assessment
should read merchant and payment facts attributed to their sources, and
where those facts conflict, return an unresolved state for human review
rather than guessing. This is the property that would have caught the
toothpaste case: payment captured, fulfillment absent, assessment not
successful.

A caveat this structure cannot escape: source attribution on evidence is
only as good as the adapter that produced it. Labeling a record as
merchant-originated does not authenticate the merchant. Evidence
authenticity is a deployment obligation, and a system that treats a
source label as proof has moved the trust problem rather than solved it.

\subsubsection{7.6 Recovery is a workflow}\label{recovery-is-a-workflow}

Prevention is cheap and reliable: a pure function evaluated before money
moves. Recovery is expensive and unreliable: it depends on merchants,
networks, and deadlines outside the system's control. The design
consequence is that effort belongs in prevention, and that recovery must
never be described as an undo.

Concretely, the structure requires that a refund \emph{request} and a
\emph{posted} refund be distinct recorded facts; that a posted refund
restore spending authority only where an approved policy says it does,
so that repeated purchase-and-refund cycles cannot evade a cumulative
limit; that disputes require explicit human authorization rather than
being filed automatically; and that a receipt report an
\emph{unrecovered} amount as a first-class figure. A system claiming
complete recovery would be misrepresenting what it controls. Reporting
the residue is the honest alternative.

\subsection{8. Fitting existing
infrastructure}\label{fitting-existing-infrastructure}

This structure is an authorization and accountability layer. It is not a
payment rail, and implementing it does not require custody, card
issuance, a new token, a blockchain, or a legal identity standard. It
sits above a rail and calls it.

That placement is deliberate, and it is what makes the structure
complementary to current agentic payment platforms rather than
competitive with them. A platform that holds funds, enforces permissions
server-side, and moves money on instruction supplies exactly the
execution layer this structure assumes. What the structure adds is
upstream of the payment call and downstream of it: an exact approved
record that determines whether the call should be made at all, and an
outcome assessment that determines what happens afterward.

Mapping onto published work in the space:

\begin{itemize}
\tightlist
\item
  \textbf{Agent transaction protocols from the card networks} are
  converging on cryptographically binding a transaction to a
  human-approved mandate specifying merchant categories, caps, and time
  windows. {[}6{]} {[}5{]} The structure here is the same idea carried
  further along the axis this paper argues matters: from categories and
  caps to the exact purchase, and from authorization to outcome.
\item
  \textbf{Rich authorization request formats} {[}2{]} and \textbf{intent
  protocols} {[}3{]} provide transport for structured authorization
  detail. An integration would require an agreed constraint schema and a
  verification algorithm, not simply an additional field.
\item
  \textbf{Workload identity and attestation work} {[}10{]} {[}11{]}
  {[}12{]} provides the machinery that would raise version binding above
  the declared level, with the freshness and coverage limits those
  specifications state.
\item
  \textbf{Payment adapters} must carry provider-specific rules for
  idempotency retention, authorization expiry, partial capture, refund
  states, and disputes. {[}16{]} {[}17{]} Local deduplication does not
  produce an exactly-once guarantee at a provider, and any field that
  cannot be faithfully enforced on a given rail should produce a
  restricted mode or a refusal rather than a claim of equivalent
  protection.
\end{itemize}

A narrowly scoped integration experiment is the appropriate next step,
and the useful output of one would be negative results: which constraint
types cannot be enforced on a given rail, what an approval surface costs
in completion rate, and where outcome evidence is unavailable or
unauthenticated. Access, provider agreement, protocol compatibility, and
commercial feasibility are all unverified here.

\subsection{9. What the prototype
establishes}\label{what-the-prototype-establishes}

A local TypeScript prototype (Scrip) implements the structure for one
purchase category --- an exact, human-approved hotel booking --- against
simulated payment and execution providers. It has never moved real
money, issued a real card, or filed a dispute. The observations below
come from source inspection and local tests.

Implemented: an approved contract canonicalized and hashed, with
approval binding to that hash; typed constraint predicates evaluated by
a pure function before any reservation; unresolved constraints blocking
execution; agent lineage, immutable version records with manifest
digests, and mandates carrying an explicit authorized-version
allow-list; authority checked before reservation and again before
dispatch; exposure reserved durably before any provider call; outcome
assessed from merchant and payment facts with the agent's narrative
excluded; refund request and posted refund recorded separately; and a
receipt reporting authorized, reserved, captured, reversed, refunded,
returned, and unrecovered amounts as distinct figures.

\subsubsection{9.1 Adversarial results}\label{adversarial-results}

An earlier revision of this code failed four attacks that a reviewer
identified by inspection. Each was reproduced as an executing test that
succeeded against the old code, then fixed. They are retained as
regression tests.

{\def\LTcaptype{none} % do not increment counter
\begin{longtable}[]{@{}
  >{\raggedright\arraybackslash}p{(\linewidth - 6\tabcolsep) * \real{0.2500}}
  >{\raggedright\arraybackslash}p{(\linewidth - 6\tabcolsep) * \real{0.2500}}
  >{\raggedright\arraybackslash}p{(\linewidth - 6\tabcolsep) * \real{0.2500}}
  >{\raggedright\arraybackslash}p{(\linewidth - 6\tabcolsep) * \real{0.2500}}@{}}
\toprule\noalign{}
\begin{minipage}[b]{\linewidth}\raggedright
Attack
\end{minipage} & \begin{minipage}[b]{\linewidth}\raggedright
Earlier behavior
\end{minipage} & \begin{minipage}[b]{\linewidth}\raggedright
Current behavior
\end{minipage} & \begin{minipage}[b]{\linewidth}\raggedright
Control
\end{minipage} \\
\midrule\noalign{}
\endhead
\bottomrule\noalign{}
\endlastfoot
Caller supplies a fabricated agent identity object & Completed a
purchase & Rejected & Service authenticates a presented credential
itself; identity type is unforgeable outside the registry \\
Credential revoked after authentication, then used & Completed a
purchase & Rejected & Credential re-verified against stored state inside
every authority check \\
Mandate scoped to refunds used to purchase & Completed a purchase &
Rejected & Requested operation and funding source checked against the
mandate \\
Operator self-asserts a stronger attestation level & Recorded as
asserted & Not expressible & Level is no longer a caller-supplied
parameter \\
\end{longtable}
}

The first of these is the instructive one. The prototype had been
reviewed as implementing version-bound authority, and it did --- but its
service boundary accepted an unauthenticated claim of having been
authenticated, so every check downstream operated on an attacker-chosen
identity. A correct authorization model behind an uncontrolled boundary
provides no protection. This is offered as a caution rather than a
solved problem: boundaries of this kind are easy to describe correctly
and easy to implement incorrectly.

\subsubsection{9.2 Limits}\label{limits}

\begin{itemize}
\tightlist
\item
  \textbf{No measured-runtime guarantee.} Registration is self-declared
  and that is the only level the registry can issue; the interface for
  stronger attestation exists with no implementation, so a stronger
  label cannot be produced. A manifest digest detects alteration of the
  record, not of the running program. An unchanged digest cannot detect
  undeclared changes, mutable model aliases, or evolving retrieved
  content.
\item
  \textbf{Human authentication is out of scope.} The principal is a
  caller-supplied identifier. Approval and revocation assume an
  authenticated surrounding service that does not exist here.
\item
  \textbf{Session-level, not request-level, credential binding.} A
  production design should bind each request to mandate, operation,
  payload digest, audience, expiry, and a replay guard.
\item
  \textbf{Evidence source labels are not authenticated.} See 7.5.
\item
  \textbf{No atomic revocation ordering with a provider.} See 7.4.
\item
  \textbf{Single category, simulated providers, no live result.} The
  natural-language-to-typed-constraint extraction described in 7.3 is
  not implemented; constraints are authored directly in fixtures.
\item
  \textbf{Canonicalization is custom.} No conformance to a published
  canonicalization scheme {[}15{]} is claimed.
\end{itemize}

\subsubsection{9.3 Local verification
record}\label{local-verification-record}

On September 18, 2026, \texttt{npm\ run\ build},
\texttt{npx\ tsc\ -\/-noEmit}, and \texttt{npm\ test} each exited 0,
with 252 tests passing and 8 skipped across 21 files and one skipped
file. The identity suite contained 26 tests, the purchase-mission suite
50; the skipped file is a PostgreSQL store suite requiring a database.

Those totals describe a test suite, not independent validations of this
paper. Passing fixtures against simulated providers do not establish
security, payment correctness, evidence authenticity, or fulfillment
reliability. No financial operation was initiated.

\subsection{10. Evaluation required}\label{evaluation-required}

The central hypothesis --- that binding authority to an exact approved
record and a declared configuration prevents failures that amount-based
controls do not --- is untested outside local fixtures. It would be
weakened if exact approval provides little protection beyond a
well-designed mandate system, or if its approval burden outweighs the
protection.

\textbf{Comparative study.} The same task set under three
configurations: amount-based limits alone; limits plus exact-record
binding; and both plus a measured deployment. Report blocked
unauthorized operations \emph{and} legitimate operations incorrectly
blocked. A control that stops every bad purchase by stopping most good
ones is not useful.

\textbf{Authority tests.} Undeclared configuration changes, mutable
aliases, stolen and expired credentials, fabricated caller identities,
operator changes, version rollback, missing scopes, and revocation
arriving during capability issuance and during dispatch.

\textbf{Failure injection.} Interrupt each persistence and network
boundary; reorder and duplicate provider events; reuse idempotency keys
past provider retention; simulate partial capture, provider
inconsistency, and refund failure. Report unresolved exposure and time
to reconcile rather than assuming timeouts are safely retryable.

\textbf{Evidence tests.} Fabricated merchant labels, compromised
adapters, conflicting confirmations, cancelled bookings, stale refund
events. Report false success, false failure, and unresolved outcomes
against independently established ground truth.

\textbf{Approval-burden study.} Measure whether people understand what
they authorized, why a change prompt appeared, and the difference
between a refund requested and money returned. Report approval counts,
completion time, abandonment, and misunderstanding alongside safety
outcomes. Given the pass\^{}k evidence in Section 3, the interesting
comparison is how burden scales with the number of purchases in a
session.

\subsection{11. Conclusion}\label{conclusion}

The public record of agent purchase failures is short, and the reasons
are structural rather than reassuring. The measured record is more
informative: agents that are usually right are considerably less often
reliably right, and the retail error category that benchmarks
over-represent is the handling of confirmation itself.

Against that, the controls currently shipped bound the amount of a
transaction. Every failure this survey located was correctly priced. The
proposal is to bind the authorization to an exact record a human
approved, evaluate it with pure functions, fail closed on anything that
cannot be represented, bind authority to a declared software
configuration, assess outcomes from evidence the agent did not author,
and treat recovery as a workflow that reports what it could not recover.

A local prototype implements this for one category against simulated
providers, and is stronger for having failed four attacks and been
fixed. It does not establish that the declared software ran, that
evidence is authentic, or that any of this is safe in production. The
next contribution should be an authenticated integration against a real
rail, and the negative results it produces.

\subsection{References}\label{references}

Primary sources below were accessed September 17-18, 2026. Repository
branches and product documentation can change; versions are stated where
visible. Source descriptions are not endorsements or independently
audited product claims.

{[}1{]} Kahlil Lalji / Natural.
\href{https://www.natural.com/blog/agentic-payments-revisited}{Agentic
payments, revisited}. August 12, 2026. Founder essay; identity,
liability, and platform sections.

{[}2{]} T. Lodderstedt, J. Richer, and B. Campbell.
\href{https://www.rfc-editor.org/rfc/rfc9396}{RFC 9396: OAuth 2.0 Rich
Authorization Requests}. 2023. Standards-track specification.

{[}3{]} Google Agentic Commerce.
\href{https://github.com/google-agentic-commerce/AP2/blob/main/docs/ap2/specification.md}{AP2
specification}. Inspected document labels itself v0.2; moving repository
branch.

{[}4{]} Google Cloud.
\href{https://cloud.google.com/blog/products/ai-machine-learning/announcing-agents-to-payments-ap2-protocol}{Announcing
Agent Payments Protocol}. Launch announcement; consulted for historical
terminology, not as the current schema.

{[}5{]} Mastercard / agent-intent contributors.
\href{https://github.com/agent-intent/verifiable-intent/}{Verifiable
Intent}. Draft v0.1, as labeled in the inspected repository.

{[}6{]} Visa.
\href{https://developer.visa.com/capabilities/trusted-agent-protocol/trusted-agent-protocol-specifications}{Trusted
Agent Protocol specifications}. Merchant-facing protocol documentation.

{[}7{]} Natural. \href{https://www.natural.com/identity}{Identity}.
Public product description.

{[}8{]} Natural. \href{https://www.natural.com/disputes}{Disputes}.
Public product description.

{[}9{]} Dan Hill / Stripe.
\href{https://stripe.com/blog/giving-agents-the-ability-to-pay}{Giving
agents the ability to pay}. April 29, 2026. Link wallet product
announcement.

{[}10{]} SPIFFE.
\href{https://spiffe.io/docs/latest/spire-about/spire-concepts/}{SPIRE
concepts}. Workload registration and attestation documentation.

{[}11{]} SLSA. \href{https://slsa.dev/spec/v1.2/provenance}{Provenance,
specification v1.2}. Software supply-chain provenance.

{[}12{]} H. Birkholz et
al.~\href{https://www.rfc-editor.org/rfc/rfc9334}{RFC 9334: Remote
ATtestation procedureS Architecture}. 2023. Informational architecture;
especially trust and freshness.

{[}13{]} M. De Rossi et
al.~\href{https://eips.ethereum.org/EIPS/eip-8004}{ERC-8004: Trustless
Agents}. Draft, created August 13, 2025; status checked at review time.

{[}14{]} B. Mohammadkhani, A. Khekade, and R. Kakkad.
\href{https://arxiv.org/html/2609.02208v1}{Agentic Settlement Protocol:
An Application Profile for Refundable, Delayed-Fulfilment Agent Commerce
on Stablecoin Rails}. arXiv:2609.02208v1, September 2, 2026.
Preprint/design proposal; no peer-review claim made here.

{[}15{]} A. Rundgren, B. Jordan, and S. Erdtman.
\href{https://www.rfc-editor.org/rfc/rfc8785}{RFC 8785: JSON
Canonicalization Scheme}. 2020. Informational specification.

{[}16{]} Stripe.
\href{https://docs.stripe.com/api/idempotent_requests}{Idempotent
requests}. API documentation; key retention and parameter matching.

{[}17{]} Stripe. \href{https://docs.stripe.com/refunds}{Refunds}.
Operational documentation; pending and failed refund states.

{[}17{]} S. Yao et
al.~\href{https://arxiv.org/abs/2406.12045}{\(\tau\)-bench: A Benchmark
for Tool-Agent-User Interaction in Real-Domain Dialogue}.
arXiv:2406.12045. Source of the pass\^{}k consistency metric and the
retail error taxonomy cited in Section 3.

{[}18{]} G. A. Fowler.
\href{https://www.washingtonpost.com/technology/2025/02/07/openai-operator-ai-agent-chatgpt/}{OpenAI's
Operator agent}. Washington Post, February 7, 2025. Firsthand account of
an unauthorized \$31.43 purchase; indexed as AI Incident Database
incident 1028.

{[}19{]} \emph{Moffatt v. Air Canada}, British Columbia Civil Resolution
Tribunal, February 2024.
\href{https://www.cbc.ca/news/canada/british-columbia/air-canada-chatbot-lawsuit-1.7116416}{CBC
report}. Tribunal ruling rejecting the ``separate legal entity''
defense.

{[}20{]} A. Smith.
\href{https://www.modernretail.co/technology/brands-are-upset-that-buy-for-me-is-featuring-their-products-on-amazon-without-permission/}{Brands
are upset that Buy for Me is featuring their products without
permission}. Modern Retail, January 6, 2026. Named merchants;
corroborated by contemporaneous CNBC reporting.

{[}21{]} M. Zeff.
\href{https://techcrunch.com/2024/12/02/the-race-is-on-to-make-ai-agents-do-your-online-shopping-for-you}{The
race is on to make AI agents do your online shopping}. TechCrunch,
December 2, 2024. Firsthand test; payment taken without fulfillment.
Single source.

{[}22{]}
\href{https://www.theregister.com/2025/07/21/replit_saastr_vibe_coding_incident/}{Replit
AI agent production database incident}. The Register, July 21, 2025.
Non-purchase; included for the irreversible-action mechanism. AI
Incident Database incident 1152.

{[}23{]} Palo Alto Networks Unit 42.
\href{https://unit42.paloaltonetworks.com/ai-agent-prompt-injection/}{Prompt
injection targeting AI agents}. March 3, 2026. Payload telemetry; the
authors state no confirmed successful real-world exploitation.

{[}24{]}
\href{https://fortune.com/2026/09/17/openai-dicloses-six-incidents-agents-going-rogue-transparency/}{OpenAI
discloses six agent incidents}. Fortune, September 17, 2026. Cited as a
negative result: no purchase or payment events among them.

{[}25{]} Brave.
\href{https://brave.com/blog/comet-prompt-injection/}{Agentic browser
security: indirect prompt injection} and
\href{https://brave.com/blog/unseeable-prompt-injections/}{unseeable
prompt injections}. 2025. Vulnerability research by a competitor to the
product examined; disclosed as such.

{[}26{]} Instinct (Spear Street Technology, Inc.).
\href{https://instinct.com/terms}{Terms of Service}. Quoted for the
appointment-as-agent, safeguard-disclaimer, non-refundable, and
liability-cap clauses in Section 4.

{[}27{]} Assistant Benchmark.
\href{https://assistantbenchmark.com/agents/instinct}{Instinct
assessment}. September 1, 2026. Cited as counter-evidence: a purchase
task passed with approval before charge.

\subsection{Appendix: prototype evidence
boundary}\label{appendix-prototype-evidence-boundary}

The local verification used the Scrip working checkout (project
directory: spending).

Commands executed, all exit status 0:

\begin{verbatim}
npm run build
npx tsc --noEmit
npm test
\end{verbatim}

Result: 252 passed, 8 skipped, across 21 passing test files and one
skipped file.

At the time of this revision the mission and identity sources remained
untracked in the working checkout, so no commit identifier designates
the tested code. This is a reproducibility gap, stated rather than
papered over: the digests below identify the inspected files, and
external replication additionally requires a published source snapshot,
dependency lockfile, runtime version, fixtures, and commands.

Source paths are relative to that checkout.

\begin{itemize}
\tightlist
\item
  Agent identity, manifest classification, attestation levels:
  \textbf{src/missions/agent-identity.ts}.
\item
  Authentication and authority: \textbf{SqliteAgentRegistry} in
  \textbf{src/missions/agent-registry.ts}, especially
  \textbf{registerVersion()}, \textbf{authenticate()},
  \textbf{verifyCredential()}, and \textbf{authorize()}.
\item
  Orchestration: \textbf{PurchaseMissionService} in
  \textbf{src/missions/purchase-mission-service.ts} --- especially
  \textbf{authenticate()}, \textbf{checkAuthority()},
  \textbf{execute()}, and \textbf{reconcile()}.
\item
  Assessment and constraint evaluation: \textbf{preflight()} and
  \textbf{assessOutcome()} in \textbf{src/missions/outcome-assessor.ts}.
\item
  Adversarial and identity fixtures:
  \textbf{tests/agent-identity.test.ts} (26 tests, including the four
  attacks in Section 9.1).
\item
  Purchase fixtures: \textbf{tests/purchase-mission-service.test.ts} (50
  tests).
\end{itemize}

SHA-256 digests of four inspected files, split across two lines for
print; concatenate without whitespace.

\begin{verbatim}
agent-identity.ts
fffaa3cba434782ed1b25bc382353ffa
d2aeee6c0dd0f9716f7f9e294aaf712f

agent-registry.ts
05ccc2b61a001e587aa26199e7a41731
98e23da532d20f34184dbba72fb51552

purchase-mission-service.ts
791405f8d8fa572b1463022012ad1189
d865934f0cffc8e790043e9323629f54

outcome-assessor.ts
9217343b5538517b74495ac5c8542222
35f1b70fb7c5028643530c5b80c78269
\end{verbatim}

\end{document}
